% Caietul de practică de specialitate. Se scrie aici și în capitole/; macheta e în caietpractica.cls (nu se modifică).
% Compilare: XeLaTeX (latexmk, Overleaf cu Compiler XeLaTeX sau automat în GitHub, la fiecare push și la fiecare etichetă).
\documentclass{caietpractica}

\titlutema{Titlul temei}
\student{Nume Prenume}
\program{Informatică}                 % sau: Matematică informatică
\traseu{a) în departament}           % sau: b) la partener: <numele partenerului>; c) în școală: <numele școlii>
\tutore{}                             % doar la traseele b) și c): Nume Prenume, funcția
\semestrul{1}                         % 2 în semestrul al doilea (Practica de specialitate 2)
\InputIfFileExists{versiune.tex}{}{}  % la o etichetă (S1-E3), GitHub scrie aici versiunea; nu se completează de mână

\begin{document}
\maketitle
\tableofcontents

\chapter*{Declarație}
\addcontentsline{toc}{chapter}{Declarație}
\markboth{Declarație}{}
Declar că acest caiet de practică și rezultatele prezentate sunt realizate de mine, sub coordonarea
conf.~univ.~dr.~Cosmin Nicolae Sabo, că toate sursele utilizate sunt citate și că am declarat mai jos instrumentele
de asistență utilizate.

\textbf{Originea temei.} Tema este stabilită în urma colaborării cu firma parteneră și aleasă pentru stagiul de practică de specialitate, fiind validată de coordonator.

\textbf{Instrumente de asistență} (inclusiv cele bazate pe inteligență artificială), pentru ce le-am utilizat și cum am verificat rezultatul:
\begin{itemize}
  \item Git și GitHub: utilizate pentru versionarea codului și generarea automată a documentației.
  \item Asistent AI: utilizat pentru structurarea documentației în LaTeX și formularea etapelor de început, rezultatele fiind verificate și adaptate manual.
\end{itemize}
\chapter{Introducere}
\etapa{S1-E1}
Stagiul de practică de specialitate oferă cadrul necesar pentru aplicarea cunoștințelor teoretice acumulate în cadrul programului de studiu în mediul real de dezvoltare software, fiind util atât pentru dezvoltarea mea profesională, cât și pentru compania parteneră prin livrarea unei soluții software funcționale.

\section{Problema}
Dezvoltarea unei aplicații software adaptate cerințelor specifice identificate în cadrul companiei, vizând optimizarea unor fluxuri de lucru sau realizarea unui produs digital nou.

\section{Obiectivele}
\begin{enumerate}
  \item Identificarea și documentarea cerințelor funcționale și nefuncționale ale temei alese.
  \item Proiectarea și implementarea arhitecturii software utilizând instrumente moderne de programare.
  \item Testarea riguroasă a componentelor realizate și integrarea corectă a acestora.
\end{enumerate}

\begin{ideeprincipala}
Dezvoltarea, implementarea și validarea unei soluții software practice care să răspundă cerințelor stabilite împreună cu tutorele de practică.
\end{ideeprincipala}
\chapter{Documentare și stadiul actual}
\etapa{S1-E2}

În cadrul acestei etape inițiale, documentarea se axează pe înțelegerea conceptelor fundamentale de Asigurare a Calității (QA) în mediul profesional, pe stadiul actual al testării software (manuale și automate) și pe stabilirea modului în care procesele moderne integrează instrumente bazate pe inteligență artificială.

\section{Surse și soluții existente}
Testarea software reprezintă un pilon esențial în dezvoltarea aplicațiilor moderne, asigurând livrarea unor produse stabile și conforme cu cerințele utilizatorilor. Conform literaturii de specialitate și bunelor practici din industrie (precum standardele ISTQB), procesul de QA combină:
\begin{itemize}
    \item Testarea manuală exploratorie și bazată pe scenarii/user stories.
    \item Testarea automatizată a interfețelor (UI) și a serviciilor web (API) pentru creșterea eficienței.
    \item Utilizarea instrumentelor moderne și a tehnicilor de prompt engineering asistate de AI pentru generarea datelor de test și optimizarea cazurilor de testare.
\end{itemize}

\section{Planul individual de lucru}
Planul de practică se desfășoară pe parcursul celor două semestre (112 ore total, împărțite în 14 săptămâni pe semestru) și este structurat pe următoarele etape:
\begin{itemize}
    \item \textbf{E1:} Tema, traseul, depozitul de lucru, convenția și stabilirea regulilor de acces.
    \item \textbf{E2:} Documentarea, fundamentarea testării și explorarea inițială a aplicației.
    \item \textbf{E3:} Analiza cerințelor și proiectarea soluției (șabloane de teste și raportarea defectelor).
    \item \textbf{E4:} Primele componente funcționale (scrierea seturilor de teste și execuția lor).
    \item \textbf{E5:} Integrarea, datele de test și structura caietului.
    \item \textbf{E6:} Versiunea de etapă testată și testele de regresie.
    \item \textbf{E7:} Caietul semestrului și pregătirea prezentării (demo).
\end{itemize}
\chapter{Analiza problemei}
\etapa{S1-E3}
Cerințele, datele, constrângerile și criteriile după care se evaluează soluția.

\chapter{Proiectarea soluției}
\etapa{S1-E3}
Modelul, arhitectura sau proiectul didactic; alegerile făcute și justificarea lor.

% Figurile se adaugă în caiet/figuri/:
% \begin{figure}[ht]\centering
%   \includegraphics[width=.8\linewidth]{figuri/arhitectura.png}
%   \caption{Arhitectura aplicației.}\label{fig:arhitectura}
% \end{figure}

\chapter{Realizarea}
\etapa{S1-E4, S1-E5}
Ce ați implementat sau realizat, cu fragmentele relevante, figuri și tabele. Codul complet stă în \texttt{produs/};
în caiet adăugați doar partea care merită explicată.

\begin{cod}[highlightlines={2}]
def exemplu(n):
    # suma primelor n numere naturale; rândul discutat e evidențiat
    return sum(range(n))
\end{cod}

\chapter{Testarea și evaluarea}
\etapa{S1-E6}
Cum ați verificat că soluția funcționează: datele de test, rezultatele, comparațiile și limitele.

\begin{table}[ht]\centering\sffamily\small
\begin{tabular}{@{}lr@{}}
\toprule
Scenariul & Rezultatul \\
\midrule
<…> & <…> \\
\bottomrule
\end{tabular}
\caption{Rezultatele testelor.}
\end{table}

\chapter{Concluzii și continuare}
\etapa{S1-E7}
Ce ați realizat, ce ați învățat, ce rămâne de făcut și cum poate continua tema în semestrul 2
(și, eventual, în lucrarea de licență).


\bibliographystyle{plain}
\bibliography{bibliografie}

\appendix
\chapter{Activitatea pe etape}
Pentru fiecare etapă: eticheta din depozit și ce s-a realizat. Data predării este cea din Campus Virtual.

\begin{center}\sffamily\small
\begin{tabular}{@{}l p{0.75\linewidth}@{}}
\toprule
Eticheta & Ce s-a realizat \\
\midrule
S1-E1 & \\
S1-E2 & \\
S1-E3 & \\
S1-E4 & \\
S1-E5 & \\
S1-E6 & \\
S1-E7 & \\
\bottomrule
\end{tabular}
\end{center}

\end{document}
